\documentclass[a4paper]{article}

% Shared Preamble for MathLatexDocs

% --- Core Packages ---
\usepackage[english]{babel}
\usepackage[utf8]{inputenc}
\usepackage{amsmath}
\usepackage{amsthm}
\usepackage{amssymb}
\usepackage{mathtools}
\usepackage{graphicx}
\usepackage{microtype} % Better justification and fewer overfull hboxes
\usepackage{bm}
\usepackage{enumitem}
\usepackage{booktabs}
\usepackage{mathrsfs}

% --- Layout & Design ---
\usepackage{docmute}
\usepackage{subfiles} % Support for master/subfile structure
\usepackage[colorinlistoftodos]{todonotes}
\usepackage[colorlinks=true, linkcolor=blue, citecolor=magenta]{hyperref}
\usepackage{cleveref} % Better cross-referencing
\usepackage{tcolorbox} % For styled boxes
\usepackage{tikz} % For drawing diagrams
\usetikzlibrary{arrows.meta,positioning}

% --- Theorem Environments (Classic Style) ---
\makeatletter
\@ifundefined{classicthm}{%
  \newtheorem{classicthm}{Theorem}[section]
  \newtheorem{classicdef}[classicthm]{Definition}
  \newtheorem{classiclem}[classicthm]{Lemma}
  \newtheorem{theorem}{Theorem}[section]
  \newtheorem{corollary}{Corollary}[theorem]
  \newtheorem{lemma}[theorem]{Lemma}
  \newtheorem{definition}{Definition}[section]
  \newtheorem{proposition}[theorem]{Proposition}
  \newtheorem{remark}[theorem]{Remark}
  \newtheorem{example}[theorem]{Example}
}{}

% Wrappers to maintain compatibility with the {Title}{Label} syntax used in files
\@ifundefined{mytheorem}{%
  \newenvironment{mytheorem}[2]{%
    \begin{classicthm}[#1]\label{thm:#2}
  }{%
    \end{classicthm}
  }
  
  \newenvironment{mydefinition}[2]{%
    \begin{classicdef}[#1]\label{def:#2}
  }{%
    \end{classicdef}
  }
  
  \newenvironment{mylemma}[2]{%
    \begin{classiclem}[#1]\label{lem:#2}
  }{%
    \end{classiclem}
  }
}{}
\makeatother

% --- Custom Commands ---
\providecommand{\Z}{\mathbb{Z}}
\providecommand{\R}{\mathbb{R}}
\providecommand{\C}{\mathbb{C}}
\providecommand{\N}{\mathbb{N}}
\providecommand{\G}{\mathcal{G}}
\providecommand{\Lop}{\mathcal{L}}
\providecommand{\PP}{\mathcal{P}}
\providecommand{\dd}{\,\mathrm{d}}
\providecommand{\bitand}{\mathbin{\&}}
\providecommand{\CFDots}{\genfrac{}{}{0pt}{0}{}{\ddots}}


\title{Symmetric Positive-Definite matrix bounds}
\author{Emil Kerimov}
\date{\today}
\begin{document}
\maketitle



\section{Motivating Problem: Correlation Bounds}
Suppose there are three assets, $a,b,c$ , by knowing the volatility of each of the assets, and the correlation between assets $a$ and $b$, as well as the correlation between assets $b$ and $c$, what is the interval of valid correlations between and assets $a$ and $c$?

This document presents one general method of addressing this problem, and solves the case of three assets analytically.

\section{Definitions and Positive-Definite Properties}

\begin{definition}\label{spd:PD}
A real symmetric $n \times n$ matrix $M$ is \textbf{positive-definite} (denoted $M \succ 0$) if for every non-zero vector $x \in \R^n$,
\begin{equation}
x^T M x > 0.
\end{equation}
\end{definition}

\begin{definition}\label{spd:PSD}
A real symmetric $n \times n$ matrix $M$ is \textbf{positive-semidefinite} (or semi-positive-definite, denoted $M \succeq 0$) if for every vector $x \in \R^n$,
\begin{equation}
x^T M x \geq 0.
\end{equation}
\end{definition}

\begin{definition}\label{spd:principle minor}
For an $n \times n$ matrix $M$, a \textbf{principal submatrix} of order $k$ (where $1 \leq k \leq n$) is obtained by choosing an index subset $I \subseteq \{1, 2, \dots, n\}$ with $|I| = k$ and keeping the entries $(M_{i,j})_{i,j \in I}$ (deleting the $n-k$ rows and columns not in $I$). The determinant of this submatrix is called a \textbf{principal minor} of order $k$.

When $I = \{1, 2, \dots, k\}$, the resulting determinant is called a \textbf{leading principal minor} of order $k$.
\end{definition}

\section{Background Matrix Theory}
We make use of some known properties of positive-definite (PD) and positive-semi-definite (PSD) matrices. 

\begin{theorem} \label{spd:PD implies positive eigenvalues} 
A PD matrix must have all positive eigenvalues.

Proof:
\begin{gather*}
\shortintertext{Let $\lambda$ be an eigenvalue of the symmetric matrix $M$ with corresponding non-zero eigenvector $v \in \R^n$:}
M v = \lambda v
\\
\shortintertext{Multiplying on the left by $v^T$:}
v^T M v = v^T (\lambda v) = \lambda (v^T v)
\\
\shortintertext{From the definition of PD, since $v \neq 0$, we have $v^T M v > 0$. Since $v \neq 0$, the Euclidean inner product $v^T v = \|v\|^2 > 0$:}
\lambda = \frac{v^T M v}{v^T v} > 0
\end{gather*}
Hence, all eigenvalues of a positive-definite matrix are strictly positive. (For a PSD matrix, $v^T M v \geq 0$ similarly implies $\lambda \geq 0$.)
\end{theorem}

\begin{theorem} \label{spd:PD implies positive principle minors} 
A PD matrix must have all positive principle minors.

Proof:
\begin{gather*}
\shortintertext{Let $M$ be an $n \times n$ PD matrix, and let $M_I$ be a $k \times k$ principal submatrix corresponding to index set $I \subseteq \{1, \dots, n\}$ with $|I| = k$. For any non-zero vector $y \in \R^k$, construct $x \in \R^n$ by placing the entries of $y$ into the positions indexed by $I$ and zeros elsewhere. Since $y \neq 0$, we have $x \neq 0$:}
y^T M_I y = x^T M x > 0
\\
\shortintertext{Because $y^T M_I y > 0$ for every non-zero $y \in \R^k$, the principal submatrix $M_I$ is itself positive-definite. By Theorem~\ref{spd:PD implies positive eigenvalues}, all $k$ eigenvalues $\lambda_1, \dots, \lambda_k$ of $M_I$ are strictly positive. Since the determinant is the product of eigenvalues:}
\det(M_I) = \prod_{i=1}^{k} \lambda_i > 0
\end{gather*}
Thus, all principal minors of a positive-definite matrix are strictly positive. (Similarly, for a PSD matrix, all principal minors are non-negative: $\det(M_I) \geq 0$.)
\end{theorem}

\section{Principal Minors Method}\label{spd:sec:main-idea}
All principle majors of a PD matrix, must be positive.
Therefore, compute them all, and figure out bounds on the unknown quantities.

\subsection{Two Assets (\texorpdfstring{$2 \times 2$}{2x2} Case)}\label{spd:subsec:two-asset}
The case $n=2$ is trivial, due to the definition of correlation being bounded by $-1$ and $1$.

\begin{gather*}
|\rho_{a,b}| \leq 1 \\
|\frac{\sigma_{a,b}}{\sigma_a \cdot \sigma_b}| \leq 1 \\
\shortintertext{Since $\sigma_a$ and $\sigma_b$ are positive by definition}
|\sigma_{a,b}| \leq \sigma_a \cdot \sigma_b
\end{gather*}

We show that the positive principle minors method arrives at the same bound. 

Let matrix $C$ be a $2$x$2$ matrix, 
\begin{equation}
C = \begin{pmatrix}
\sigma_a^2 & \sigma_{ab} \\
\sigma_{ab} & \sigma_b^2
\end{pmatrix}\label{spd:eq:equation}
\end{equation}

For $C$ to be positive-semidefinite, its determinant (a principal minor) must be non-negative:
\begin{gather*}
\det(C) = \sigma_a^2 \sigma_b^2 - \sigma_{ab}^2 \geq 0 \\
\sigma_{ab}^2 \leq \sigma_a^2 \sigma_b^2 \\
|\sigma_{ab}| \leq \sigma_a \cdot \sigma_b
\end{gather*}

\subsection{Three Assets (\texorpdfstring{$3 \times 3$}{3x3} Case)}
We can write the correlation matrix of three assets as:

\begin{equation}
C = \begin{pmatrix}
1 & \rho_{ab} & \mathbf{\rho_{ac}} \\
\rho_{ab} & 1 & \rho_{bc} \\
\mathbf{\rho_{ac}} & \rho_{bc} & 1
\end{pmatrix}
\end{equation}

\begin{theorem} \label{spd:three} 
\begin{equation}
\boxed{
|\mathbf{\rho_{ac}} - \rho_{ab} \cdot \rho_{bc}| \leq \sqrt{(1 - \rho_{ab}^2) \cdot (1 - \rho_{bc}^2)}
}
\end{equation}

Proof:
\begin{gather*}
\shortintertext{For the correlation matrix $C$ to be positive-semidefinite, its determinant must be non-negative:}
\det(C) \geq 0
\\
1 \cdot (1 - \rho_{bc}^2) - \rho_{ab} (\rho_{ab} - \mathbf{\rho_{ac}} \rho_{bc}) + \mathbf{\rho_{ac}} (\rho_{ab} \rho_{bc} - \mathbf{\rho_{ac}}) \geq 0
\\
1 - \rho_{bc}^2 - \rho_{ab}^2 + 2 \rho_{ab} \rho_{bc} \mathbf{\rho_{ac}} - \mathbf{\rho_{ac}}^2 \geq 0
\\
\shortintertext{Rearranging in terms of the unknown $\mathbf{\rho_{ac}}$:}
-\mathbf{\rho_{ac}}^2 + 2(\rho_{ab} \rho_{bc}) \mathbf{\rho_{ac}} + (1 - \rho_{ab}^2 - \rho_{bc}^2) \geq 0
\\
\mathbf{\rho_{ac}}^2 - 2(\rho_{ab} \rho_{bc}) \mathbf{\rho_{ac}} \leq 1 - \rho_{ab}^2 - \rho_{bc}^2
\\
\shortintertext{Completing the square:}
(\mathbf{\rho_{ac}} - \rho_{ab} \rho_{bc})^2 \leq 1 - \rho_{ab}^2 - \rho_{bc}^2 + (\rho_{ab} \rho_{bc})^2
\\
\shortintertext{Factoring the right-hand side:}
1 - \rho_{ab}^2 - \rho_{bc}^2 (1 - \rho_{ab}^2) = (1 - \rho_{ab}^2)(1 - \rho_{bc}^2)
\\
(\mathbf{\rho_{ac}} - \rho_{ab} \rho_{bc})^2 \leq (1 - \rho_{ab}^2)(1 - \rho_{bc}^2)
\\
\shortintertext{Since $1 - \rho_{ab}^2 \geq 0$ and $1 - \rho_{bc}^2 \geq 0$ from the $2 \times 2$ principal minors, taking the square root gives:}
|\mathbf{\rho_{ac}} - \rho_{ab} \cdot \rho_{bc}| \leq \sqrt{(1 - \rho_{ab}^2) \cdot (1 - \rho_{bc}^2)}
\end{gather*}
\end{theorem}

\subsection{Four Assets (\texorpdfstring{$4 \times 4$}{4x4} Case)}
We can write the correlation matrix of four assets as:  

\begin{equation}
C = \begin{pmatrix}
1 & \rho_{ab} & \rho_{ac} & x \\
\rho_{ab} & 1 & \rho_{bc} & \rho_{bd} \\
\rho_{ac} & \rho_{bc} & 1 & \rho_{cd} \\
x & \rho_{bd} & \rho_{cd} & 1
\end{pmatrix}
\end{equation}

\begin{theorem} \label{spd:four} 
Bounds on $x=\rho_{ad}$
\\
\boxed{
\begin{gathered}
% \boxed{
|\mathbf{\rho_{ad}} - \frac{(\rho_{a,b} \cdot \rho_{b,d} + \rho_{a,c} \cdot \rho_{c,d}) - \rho_{b,c} \cdot (\rho_{a,b} \cdot \rho_{c,d} + \rho_{a,c} \cdot \rho_{b,d}) }{1 - \rho_{b,c}^2} | \leq (1 - \rho_{b,c}^2) \sqrt{D}
\\
D = (1 - \rho_{a,b}^2 - \rho_{a,c}^2 -\rho_{b,c}^2 + 2 \cdot \rho_{a,b} \cdot \rho_{a,c} \cdot \rho_{b,c}) \cdot (1 - \rho_{b,c}^2 - \rho_{b,d}^2- \rho_{c,d}^2+ 2 \cdot \rho_{b,c} \cdot \rho_{b,d} \cdot \rho_{c,d})
\end{gathered}
}
   
Proof:
\begin{gather*}
Det(C) \geq 0 
\\
1 - (\rho_{c,d}^2 + \rho_{b,c}^2 + \rho_{b,d}^2 + \rho_{a,b}^2 + \rho_{a,c}^2) 
\\
+ 2 \cdot (\rho_{b,c} \rho_{b,d} \rho_{c,d} 
+ \rho_{a,b} \rho_{b,c} \rho_{a,c}) 
\\
- 2 ( \rho_{a,b} \rho_{b,d} \rho_{a,c} \rho_{c,d} )
\\
+ (\rho_{a,b}^2 \rho_{c,d}^2 + \rho_{a,c}^2 \rho_{b,d}^2)
\\
+ 2x (\rho_{a,b} \rho_{b,d} + \rho_{a,c} \rho_{c,d})
\\
- 2x ( \rho_{a,b} \rho_{b,c} \rho_{c,d} 
+ \rho_{a,c} \rho_{b,d} \rho_{b,c} )
\\
+ x^2 ( \rho_{b,c}^2 -1) \geq 0
\\
\shortintertext{This is just a quadradic equation}
\\
a = \rho_{b,c}^2 -1 
\\
b = 2( \rho_{a,b} \rho_{b,d} + \rho_{a,c} \rho_{c,d}) - 2 ( \rho_{a,b} \rho_{b,c} \rho_{c,d} 
+ \rho_{a,c} \rho_{b,d} \rho_{b,c} ) 
\\
c = 1 - (\rho_{c,d}^2 + \rho_{b,c}^2 + \rho_{b,d}^2 + \rho_{a,b}^2 + \rho_{a,c}^2) 
\\
+ 2 \cdot (\rho_{b,c} \rho_{b,d} \rho_{c,d} 
+ \rho_{a,b} \rho_{b,c} \rho_{a,c}) 
\\
- 2 ( \rho_{a,b} \rho_{b,d} \rho_{a,c} \rho_{c,d} )
\\
+ (\rho_{a,b}^2 \rho_{c,d}^2 + \rho_{a,c}^2 \rho_{b,d}^2)
\\
ax^2 + bx + c \geq 0
\\
|2ax + b| \leq \sqrt{b^2 - 4ac}
\\
\shortintertext{Dividing all sides by 2}
a = \rho_{b,c}^2 -1 
\\
b = ( \rho_{a,b} \rho_{b,d} + \rho_{a,c} \rho_{c,d}) - ( \rho_{a,b} \rho_{b,c} \rho_{c,d} 
+ \rho_{a,c} \rho_{b,d} \rho_{b,c} ) 
\\
c = 1 - (\rho_{c,d}^2 + \rho_{b,c}^2 + \rho_{b,d}^2 + \rho_{a,b}^2 + \rho_{a,c}^2) 
\\
+ 2 \cdot (\rho_{b,c} \rho_{b,d} \rho_{c,d} 
+ \rho_{a,b} \rho_{b,c} \rho_{a,c}) 
\\
- 2 ( \rho_{a,b} \rho_{b,d} \rho_{a,c} \rho_{c,d} )
\\
+ (\rho_{a,b}^2 \rho_{c,d}^2 + \rho_{a,c}^2 \rho_{b,d}^2)
\\
|ax + b| \leq \sqrt{b^2 - ac}
\\
\shortintertext{Simplifying the square root}
b^2 - ac = (a,b)^2 (b,d)^2 + (a,c)^2 (c,d)^2 + (b,c)^2 (c,d)^2 +(b,c)^2 (b,d)^2
\\
 + (b,c)^2 (a,b)^2 + (b,c)^2 (a,c)^2 - (a,b)^2 (c,d)^2 -(a,c)^2 (b,d)^2 
\\
+(2*(b,c)^2 + (c,d)^2 + (b,d)^2 + (a,b)^2 + (a,c)^2) -1 
\\
+
4 (a,b)(b,c)^2 (c,d)(a,c)(b,c) - 2((b,c)(b,d)(c,d) + (a,b)(b,c)(a,c)) - (b,c)^4 
\\
- 2((a,b)^2 (b,d)(b,c)(c,d) + (a,b)(b,d)^2(b,c) (a,c) + (a,c)(c,d)^2 (a,b) (b,c) 
\\
+ (a,c)^2 (c,d) (b,c) (b,d)) - 2 ( (b,c)^3 (b,d) (c,d) + (b,c)^3 (a,b) (a,c))
\end{gather*}
\end{theorem}

\end{document}
